% TOP_PROBLEM: {"id": 2302075, "title": "Research Problems in Function Theory — Problem 2.75", "category_id": 8, "category": {"id": 8, "name": "analysis", "display_name": "Analysis", "description": "Limits, continuity, calculus, and function theory.", "slug": "analysis", "order_index": 8, "created_at": "2026-07-31T15:26:25.671Z"}, "status": "open"}
% FIRST_POSTED: null
% ENTRY_KIND: statement_only
% Imported from: batch12.tex; UnsolvedMath ID 2302075
% Compile twice: lualatex 2302075.tex
\documentclass[11pt]{article}
\usepackage{fontspec}
\setmainfont{Latin Modern Roman}
\usepackage{amsmath,amssymb,mathtools,unicode-math}
\setmathfont{Latin Modern Math}
\usepackage{xurl,hyperref}
\hypersetup{colorlinks=true,urlcolor=blue,pdftitle={UnsolvedMath Problem 2302075}}
\newcommand{\sref}[2]{\hyperlink{source:#1:#2}{[S#2]}}
\newcommand{\cataloguescope}[1]{\paragraph{Mathematical question.}#1}
\begin{document}
% UnsolvedMath ID 2302075; source catalogue code AMR-022-2075
\setcounter{section}{2302074}
\section{Research Problems in Function Theory — Problem 2.75}\label{2302075}
{\small\sffamily UnsolvedMath ID 2302075 · Analysis}
\subsection{Definitions and mathematical statement}

\paragraph{Shared notation.}$\mathbb N=\{1,2,\ldots\}$, and $\mathbb Z,\mathbb Q,\mathbb R,\mathbb C$ denote the integers, rationals, reals and complex numbers. Any other conventions are specified locally.


For an entire function $f$, put $M(r,f)=\max_{|z|=r}|f(z)|$, $T(r,f)=(2\pi)^{-1}\int_0^{2\pi}\log^+|f(re^{i\theta})|\,d\theta$, and $\rho(f)=\limsup_{r\to\infty}\log\log M(r,f)/\log r$, where $\log^+t=\max(0,\log t)$. Let $\rho(r)$ be a finite proximate order: it is eventually nonnegative, continuous and piecewise continuously differentiable, $\rho(r)\to\rho\in[0,\infty)$, and $r\rho'(r)\log r\to0$ wherever the derivative is defined. Say that $f$ has proximate order $\rho(r)$ when $0<\limsup_{r\to\infty}\log M(r,f)/r^{\rho(r)}<\infty$. Define its indicator by
\[h(\theta)=\limsup_{r\to\infty}\frac{\log|f(re^{i\theta})|}{r^{\rho(r)}}.\]
\paragraph{Statement recorded in the supplied source.}
Determine all indicator functions $h$ that can arise from an entire, locally uniformly convergent exponential series
\[f(z)=\sum_{n=1}^{\infty}a_n e^{\lambda_n z},\qquad a_n>0,\quad0\leq\lambda_1<\lambda_2<\cdots\to\infty,\]
of finite proximate order. \sref{2302075}{2}
\subsection{Short English statement}
Classify the possible growth indicators of entire exponential series with positive coefficients and strictly increasing nonnegative exponents.
\subsection{Sources}
\begin{description}
\item[\hypertarget{source:2302075:1}{S1}] ULAM.AI and the UnsolvedMath Contributors, UnsolvedMath: A Curated Collection of Open Mathematics Problems, record 2302075 (AMR-022-2075). CC BY 4.0. \url{https://huggingface.co/datasets/ulamai/UnsolvedMath}.
\item[\hypertarget{source:2302075:2}{S2}] W. K. Hayman and E. F. Lingham, \emph{Research Problems in Function Theory}, arXiv:1809.07200v2 (2018), Problem 2.75. \url{https://arxiv.org/abs/1809.07200}.
\end{description}
\label{end:2302075}
\end{document}
